\uuid{ixmC}
\titre{Équation droite/plan}
\auteur{bodin ; hanani ; mzari}
\organisation{exo7}
\licence{CC BY-NC-SA 4.0}
\datecreate{2020-01-20}
\niveau{L1}
\module{Géométrie}
\chapitre{Géométrie de l'espace}
\sousChapitre{Plans}
\difficulte{3}
\theme{Plans -- Droites, équation droite, plan}
\exo7qcmid{lille-1.geometrie-espace.015}

\contenu{
\texte{
Pour ces questions, l'espace est muni d'un repère orthonormé direct  $(O, \vec{i}, \vec{j}, \vec{k})$.
}

\begin{question}
\qclassification{141.03,141.04}
\qtags{motcle=équation droite/plan}
Soit $P$ et $P'$ deux plans non parallèles d'équations :  $ax+by+cz+d =0$ et $a'x+b'y+c'z+d' =0$ respectivement. Soit $D=P\cap P'$ et $Q$ un plan contenant $D$. Quelles sont les assertions vraies ?
\begin{answers}

     \bad{Une équation cartésienne de $Q$ est $ax+by+cz+d =0$.}

     \bad{Une équation cartésienne de $Q$ est $a'x+b'y+c'z+d' =0$.}

     \good{Une équation cartésienne de $Q$ est  de la forme : $\alpha(ax+by+cz+d)+\beta(a'x+b'y+c'z+d')=0,$  où $ \alpha, \beta \in \Rr$  tels que  $(\alpha a+ \beta a', \alpha b+ \beta b', \alpha c+ \beta c', \alpha d+ \beta d') \neq (0,0,0,0)$.}

      \good{Si $Q\neq P'$, une équation cartésienne de $Q$ est de la forme :
       $(ax+by+cz+d)+\alpha(a'x+b'y+c'z+d')=0, $ où $ \alpha \in \Rr$  tel que $ (a+ \alpha a',  b+ \alpha b', c+ \alpha c',  d+ \alpha d') \neq (0,0,0,0)$.}

\end{answers}
\begin{explanations}
$\vec{n}(a,b,c)$ est un vecteur normal à $P$ et $\vec{n'}(a',b',c')$ est un vecteur normal à $P'$, donc un vecteur normal à $Q$ est une combinaison linéaire de $\vec{n}$ et $\vec{n'}$. Par conséquent, une équation cartésienne de $Q$ est de la forme :  $\alpha(ax+by+cz)+\beta(a'x+b'y+c'z) + \gamma =0.$ D'autre part, si $A(x_0,y_0,z_0) \in D$, $A \in Q$. On déduit que $\gamma = \alpha d + \beta d'$.
\end{explanations}
\end{question}
}
